\section{Methods}

\subsection{Flow records and features}
A packet is stored in 20 bytes: a timestamp in microseconds relative to the
start of its batch, the two addresses of Internet Protocol version 4 (IPv4),
the two ports, the protocol, the TCP flags and the packet length. A flow is
unidirectional and is identified by the five key fields. A batch is a time
window of a capture; flows are defined inside a batch, as flow exporters do
with their timeouts~\cite{netflow}.

For every flow we keep a 176-byte record of \emph{exact integer accumulators}
and derive the \nfeat{} features of Table~\ref{tab:features} from it. A burst
is a maximal run of packets of the flow whose gaps are shorter than
1\,ms. The standard deviations are computed from integer sums with a 128-bit
numerator, with the same code on the processor and on the GPU. This design
makes results comparable bit by bit instead of ``close enough''.

\begin{table}[H]
\centering\footnotesize
\caption{Features computed for every flow. The last column tells whether the
feature depends on the order of the packets inside the flow.}\label{tab:features}
\begin{tabular}{@{}p{0.23\linewidth}p{0.55\linewidth}c@{}}
\toprule
Group & Features & Order \\
\midrule
Volume & packets, bytes, packet rate, byte rate & no \\
Time & first and last timestamp, duration & no \\
Packet length & minimum, maximum, mean, standard deviation & no \\
Flags & packets with FIN, SYN, RST, PSH, ACK, URG & no \\
Entropy & Shannon entropy of the length histogram (16 bins of 100 bytes) & no \\
Inter-arrival time & minimum, maximum, mean, standard deviation & yes \\
Bursts & number of bursts, longest burst, mean burst size & yes \\
\bottomrule
\end{tabular}
\end{table}

\subsection{Data sources}
\emph{Synthetic traffic.} A generator produces $n$ packets over $F$ flows. The
flow of each packet is drawn from a Zipf distribution: the flow of popularity
rank $r$ is chosen with probability proportional to $r^{-s}$. The exponent $s$
controls the skew: $s=0$ is uniform, and at $s=1.5$ one flow receives
\topshareHi\% of all packets. Arrivals form a Poisson process of 500{,}000
packets per second. The generator is counter-based, so a seed always gives the
same packets.

\emph{Real traffic.} We use a public 15-minute trace of a transit link of the
WIDE backbone in Japan, from the MAWI archive~\cite{mawi} (sample point F,
January 1, 2024, 14:00). It has 129.1 million packets, of which 114.4 million
are IPv4. A converter (\code{pcap2pkt}) parses the capture file and writes
batches. One minute of this trace has \realN{} million packets in \realF{}
million flows: \realSingle\% of the flows have a single packet, while the
largest flow has \realTop\% of the packets. Real traffic therefore combines
both difficulties: very many flows and elephants.

\subsection{Sequential baseline}
The baseline reads the packets in time order. For each packet it looks up the
flow in an open-addressing hash table with linear probing and updates the
record in place, including the last timestamp and the length of the current
burst. This is the reference for correctness and for speedup. All processor
versions use the same table and the same update function.

\subsection{OpenMP strategies}
A first idea is one shared table protected by locks or atomic instructions. It
cannot compute the order-dependent features, because threads would update a
flow out of order, and it contends on popular flows. We measure that contention
separately (Section~III-\ref{sec:contention}) and use two other strategies, shown in
Fig.~\ref{fig:pipeline}.

\emph{Merge} (\code{omp-merge}). The packets are split into $T$ contiguous
chunks, one per thread. Each thread aggregates its chunk into a private table.
The partial records of a flow are then combined in chunk order with a merge
operation. To make the merge exact for order-dependent features, a partial
record also stores the lengths of its first and last burst. When two partial
records are merged, the gap between the last packet of the first and the first
packet of the second gives one more inter-arrival time, and decides whether
the two boundary bursts join. The operation is associative, so it is a correct
parallel reduction. The merge itself runs in parallel: the hash space is cut
into ranges and each range is merged by one thread.

\emph{Partition} (\code{omp-part}). The packets are first scattered into
1{,}024 partitions by the top bits of the flow hash, with a parallel counting
sort that keeps the time order. All packets of a flow land in the same
partition. Each partition is then aggregated by one thread with a private
table, under a dynamic schedule with the largest partitions first. Nothing is
shared and nothing is merged, and a partition fits in the cache. Its weakness
is that one flow cannot be split among threads.

\begin{figure*}[t]
\centering
\begin{tikzpicture}[font=\footnotesize, >=Stealth,
  box/.style={draw=ink2, rounded corners=2pt, minimum height=6.5mm, inner xsep=4pt, align=center},
  lab/.style={font=\footnotesize\itshape, text=ink2, anchor=east}]
% row 1: merge
\node[lab] at (-0.2,2.4) {OpenMP merge};
\node[box, fill=b1] (m1) at (1.6,2.4) {$T$ time chunks};
\node[box, fill=b1, right=7mm of m1] (m2) {private table\\per thread};
\node[box, fill=b1, right=7mm of m2] (m3) {merge partial records\\by hash range};
\node[box, fill=y1, right=7mm of m3] (m4) {derive\\features};
\draw[->] (m1)--(m2); \draw[->] (m2)--(m3); \draw[->] (m3)--(m4);
% row 2: partition
\node[lab] at (-0.2,1.2) {OpenMP partition};
\node[box, fill=o1] (p1) at (1.6,1.2) {hash packets};
\node[box, fill=o1, right=7mm of p1] (p2) {scatter into\\1{,}024 partitions};
\node[box, fill=o1, right=7mm of p2] (p3) {private table\\per partition};
\node[box, fill=y1, right=7mm of p3] (p4) {derive\\features};
\draw[->] (p1)--(p2); \draw[->] (p2)--(p3); \draw[->] (p3)--(p4);
% row 3: GPU
\node[lab] at (-0.2,0) {CUDA};
\node[box, fill=g1] (g0) at (1.25,0) {upload};
\node[box, fill=g1, right=4mm of g0] (g1) {flow id:\\lock-free table};
\node[box, fill=g1, right=4mm of g1] (g2) {stable sort\\by flow id};
\node[box, fill=g1, right=4mm of g2] (g3) {gather,\\prefix sum};
\node[box, fill=g1, right=4mm of g3] (g4) {reduce: flow /\\atomic / chunk};
\node[box, fill=y1, right=4mm of g4] (g5) {derive\\features};
\node[box, fill=g1, right=4mm of g5] (g6) {download};
\draw[->] (g0)--(g1); \draw[->] (g1)--(g2); \draw[->] (g2)--(g3); \draw[->] (g3)--(g4); \draw[->] (g4)--(g5); \draw[->] (g5)--(g6);
\end{tikzpicture}
\caption{The three parallel pipelines. Every row starts from the same
time-ordered packet batch and ends with the same flow records.}
\label{fig:pipeline}
\end{figure*}

\subsection{CUDA pipeline}
The GPU version (Fig.~\ref{fig:pipeline}) does not update a table packet by
packet. It regroups the packets so that every feature becomes a reduction over
a contiguous segment.

\emph{1. Flow identifier.} One thread per packet inserts the key in a
lock-free open-addressing table~\cite{alcantara}. A slot is a single 32-bit
word that stores the index of the packet that claimed it, so one
compare-and-swap inserts a 13-byte key, and keys are compared through that
packet. The slot index identifies the flow exactly; no hash collisions are
tolerated.

\emph{2. Sort.} A stable radix sort from Thrust~\cite{thrust} orders the packet
indices by flow identifier. The packets of a flow become contiguous and remain
in time order.

\emph{3. Gather.} Timestamps and lengths are copied into sorted order, so the
following kernels read memory sequentially (coalesced access). A prefix sum
over the flow boundaries numbers the flows.

\emph{4. Reduce.} After sorting, the inter-arrival time of a packet is the
difference with its left neighbor, so every feature is a sum, minimum or
maximum over a segment. We implemented three kernels. \code{gpu-flow} uses one
thread per flow, which walks the packets of its flow. \code{gpu-atomic} uses
one thread per packet, which adds its contribution to the record of its flow
with atomic instructions. \code{gpu-chunk} uses one thread per chunk of $c$
consecutive sorted packets; it accumulates in registers while the flow does
not change and issues one atomic update per flow met in the chunk. Burst
lengths are obtained by compacting the positions where bursts start.

\emph{5. Derive and download.} A kernel derives the statistical features, and
records are copied back to the host.

\subsection{MPI4Py distribution}
Captures are processed as independent batches, which fits distributed
memory: a driver written with MPI4Py gives batches to processes, and each
process calls the C++ kernel through a shared library. With the \emph{static}
schedule, each process takes a contiguous block of the batch list. With the
\emph{dynamic} schedule, every process takes the next unclaimed batch from a
shared counter that is advanced with a one-sided fetch-and-add operation,
largest batches first. Process 0 collects either a summary per batch or all
the flow records (\code{Gatherv}).

\subsection{Correctness and quality metric}
The application-specific quality metric is \emph{flow-record accuracy}: the
fraction of flows whose record is identical to the reference. The reference
chain has two links. First, the sequential baseline is compared with an
independent implementation written in plain Python, which groups packets in
a dictionary and computes every feature from the full packet list of the
flow. Second, every parallel version is compared with the baseline: integer
accumulators must be identical and derived features must agree within a
relative tolerance of $10^{-9}$. The test suite covers edge cases (one
packet, one flow, all flows distinct, more threads than packets, odd thread
counts, extreme burst gaps), all GPU launch configurations, the MPI driver and
the real trace. During the benchmarks, every run also reports a checksum of
its records, and the driver stops if two versions disagree on the same input.

\subsection{Code and data availability}
The source code, build instructions, configuration files, raw results and the
scripts that regenerate every figure and table are in the open repository
\url{https://github.com/Bunnydragon96/hpc-flow-feature-extraction}. The
synthetic data are generated by the programs. The real trace is downloaded
from the MAWI archive~\cite{mawi} with the command given in the repository.
